%==============================================================================
\chapter{Definición formal del problema y metodología}
\label{cap:problema}
%==============================================================================

A partir de los fundamentos anteriores, el diseño del portafolio requiere fijar
la formulación matemática del JSSP en MiniZinc y el esquema formal para
seleccionar y secuenciar las configuraciones, junto con la metodología seguida.

%------------------------------------------------------------------------------
\section{Notación y definición del JSSP}
\label{sec:notacion}
%------------------------------------------------------------------------------

Una instancia del JSSP clásico se define por un conjunto de $n$ trabajos
$J = \{1, \dots, n\}$ y un conjunto de $m$ máquinas $M = \{0, \dots, m-1\}$.
Cada trabajo $j$ es una secuencia ordenada de $m$ operaciones
$o_{j1}, \dots, o_{jm}$, y cada operación $o_{jk}$ debe procesarse en la
máquina $\mu_{jk} \in M$ durante $p_{jk} \in \mathbb{Z}^{+}$ unidades de
tiempo. En la versión clásica cada trabajo pasa exactamente una vez por cada
máquina, de modo que hay $n \cdot m$ operaciones, condición que cumplen todas
las instancias de este trabajo. Las variables de decisión son el tiempo de
inicio $s_{jk} \in \mathbb{Z}_{\ge 0}$ de cada operación y el
\textit{makespan} $C_{\max}$, el instante en que termina la última operación.

Una operación no puede empezar hasta que termine la anterior de su trabajo
(precedencia), cada máquina procesa una sola operación a la vez y sin
interrupciones (capacidad), y el \textit{makespan} no puede ser menor que el
final de ningún trabajo \citep{pinedo2022}:
\begin{align}
  & s_{jk} + p_{jk} \le s_{j,k+1},
    && \forall j \in J,\ k = 1, \dots, m-1, \label{eq:precedencia} \\
  & s_{jk} + p_{jk} \le s_{j'k'} \;\lor\; s_{j'k'} + p_{j'k'} \le s_{jk},
    && \forall\, o_{jk} \ne o_{j'k'} : \mu_{jk} = \mu_{j'k'}, \label{eq:disyuntiva} \\
  & C_{\max} \ge s_{jm} + p_{jm}, && \forall j \in J. \label{eq:makespan}
\end{align}
El objetivo es
\begin{equation}
  \min\; C_{\max}.
  \label{eq:objetivo}
\end{equation}
Las restricciones~\eqref{eq:disyuntiva} se llaman disyuntivas porque expresan
una elección entre dos órdenes posibles, y es esa elección combinatoria la que
hace difícil el problema.

El valor óptimo $C^{*}_{\max}$ admite una cota inferior simple, porque ningún
programa puede terminar antes de que el trabajo más largo complete sus
operaciones ni antes de que la máquina más cargada procese todas las suyas,
la misma cota que \citet{taillard1993} reporta para sus instancias:
\begin{equation}
  C^{*}_{\max} \ge \max\Bigl( \max_{j \in J} \sum_{k=1}^{m} p_{jk},\;
  \max_{r \in M} \sum_{(j,k) : \mu_{jk} = r} p_{jk} \Bigr).
  \label{eq:cota_inferior}
\end{equation}
Por otro lado, ejecutar todas las operaciones una tras otra siempre es
factible, así que $C^{*}_{\max} \le H = \sum_{j,k} p_{jk}$, el horizonte
trivial.

%------------------------------------------------------------------------------
\section{Modelos en MiniZinc}
\label{sec:modelo_minizinc}
%------------------------------------------------------------------------------

El trabajo usa tres modelos de MiniZinc que traducen la misma formulación y,
por tanto, tienen las mismas soluciones óptimas. Se diferencian en la
información adicional que reciben los \textit{solvers}, dando lugar al modelo
original (\texttt{jssp.mzn}), al modelo con cotas (\texttt{jssp\_cotas.mzn}),
que usa el portafolio, y al modelo con cotas y estrategia de búsqueda
(\texttt{jssp\_\allowbreak cotas\_\allowbreak busqueda.mzn}), con el que Chuffed
forma la tercera configuración del portafolio.

\subsection{Modelo original}
\label{sec:modelo_original}

En el modelo original los tiempos de inicio y el \textit{makespan} tienen
dominio $0..H$ y cada \textit{solver} decide cómo buscar. Las instancias se
entregan como archivos \texttt{.dzn} con cuatro parámetros que especifican el
número de trabajos y de máquinas junto con las matrices \texttt{machines}
($\mu_{jk}$) y \texttt{durations} ($p_{jk}$), en las que la fila $j$ describe
las operaciones del trabajo $j$ en su orden de ejecución. La precedencia se
expresa directamente con la ecuación~\eqref{eq:precedencia}. Para la
capacidad, el modelo reúne, en cada máquina, los tiempos de inicio y las
duraciones de las $n$ operaciones que la usan e impone la restricción global
\texttt{disjunctive}, equivalente al conjunto de
disyunciones~\eqref{eq:disyuntiva}, que permite a cada \textit{solver} usar un
propagador especializado. El código completo se presenta en la
Sección~\ref{sec:modelo_cotas}, ya que el modelo con cotas solo modifica
algunas líneas.

El horizonte trivial constituye una cota muy holgada, puesto que en las 242
instancias $H$ es entre 3.6 y 18.8 veces la mejor solución conocida (11.4
veces en promedio). Con dominios tan amplios, un \textit{solver} puede aceptar
como primera solución un programa casi secuencial y avanzar muy lentamente
desde ahí. Este comportamiento, que se observó en Chuffed
(Sección~\ref{sec:ejecucion_solvers}), motivó la incorporación de cotas.
La idea de acotar el modelo surgió de ese diagnóstico y se llevó a cabo con
técnicas conocidas en \textit{scheduling}: la cota inferior de la
ecuación~\eqref{eq:cota_inferior} y el ajuste de las \textit{cabezas} y
\textit{colas} de cada operación respecto a una cota superior del
\textit{makespan} \citep{carlier1994}.

\subsection{Cotas del \textit{makespan}}
\label{sec:cotas}

Antes de resolver una instancia se calculan dos cotas de su óptimo. La cota
inferior $lb$ es la de la ecuación~\eqref{eq:cota_inferior}. La cota superior
$ub$ es el \textit{makespan} de un programa factible construido con el
algoritmo de \citet{giffler1960}, que genera un programa activo, es decir, uno
en el que ninguna operación puede adelantarse sin retrasar otra. El algoritmo
programa una operación por paso:

\begin{enumerate}
  \item entre las operaciones disponibles (la siguiente de cada trabajo), toma
    la que puede terminar más temprano, $o^{*}$, y su máquina $r^{*}$;
  \item forma el conjunto de conflicto con las operaciones disponibles que
    usan $r^{*}$ y pueden empezar antes de que termine $o^{*}$;
  \item elige una de ellas con una regla de prioridad y la programa en su
    inicio más temprano, hasta programar las $n \cdot m$ operaciones.
\end{enumerate}

\noindent Se usaron seis reglas de prioridad clásicas \citep{panwalkar1977}
(Tabla~\ref{tab:reglas}). Para cada instancia se construye un programa con
cada regla, se verifica que cumpla las restricciones de precedencia y de
capacidad y se toma como $ub$ el menor \textit{makespan}. El cálculo de ambas
cotas para las 242 instancias toma pocos segundos
(\texttt{src/02\_generar\_cotas.py}).

\begin{table}[H]
  \centering
  \small
  \singlespacing
  \caption{Reglas de prioridad usadas en el algoritmo de Giffler y Thompson.}
  \label{tab:reglas}
  \begin{tabularx}{\textwidth}{l >{\raggedright\arraybackslash}X r}
    \toprule
    \textbf{Regla} & \textbf{Elige la operación\dots} & \textbf{Veces mejor} \\
    \midrule
    SPT  & de menor duración & 0 \\
    LPT  & de mayor duración & 0 \\
    MWKR & del trabajo con más tiempo de proceso pendiente & 91 \\
    LWKR & del trabajo con menos tiempo de proceso pendiente & 3 \\
    MOR  & del trabajo con más operaciones pendientes & 58 \\
    FIFO & que quedó disponible primero & 90 \\
    \bottomrule
  \end{tabularx}
  \par\vspace{2pt}{\footnotesize Veces mejor: número de instancias en que la regla
  dio la menor $ub$ (en caso de empate se cuenta la primera regla de la lista).}
\end{table}

En promedio, la cota inferior equivale al 88.3\,\% de la mejor solución
conocida y la cota superior la supera en un 26.3\,\% (detalle por familia en
la Tabla~\ref{tab:cotas}). En seis instancias (\texttt{la06}, \texttt{la10},
\texttt{la12}, \texttt{la14}, \texttt{swv16} y \texttt{swv18}) ambas cotas
coinciden, de modo que el programa heurístico es óptimo y su optimalidad queda
probada sin ejecutar ningún \textit{solver}.

\subsection{Modelo con cotas}
\label{sec:modelo_cotas}

Un aporte realizado en este trabajo es integrar, mediante técnicas habituales
de reducción de dominios, las cotas del preprocesamiento de Giffler y Thompson
en un modelo con cotas en MiniZinc que recibe $lb$ y $ub$ como parámetros: los tiempos de inicio pasan a $0..ub$, el
\textit{makespan} queda en $lb..ub$ y cada operación no puede empezar antes de
las anteriores de su trabajo ni terminar después de $ub$ restadas las
siguientes, lo que corresponde a acotar sus cabezas y colas \citep{carlier1994},
\begin{equation}
  \sum_{l=1}^{k-1} p_{jl} \;\le\; s_{jk} \;\le\; ub - \sum_{l=k}^{m} p_{jl},
  \qquad \forall j \in J,\ k = 1, \dots, m.
  \label{eq:cabeza_cola}
\end{equation}
El código del modelo, sin la sección de salida, se presenta a continuación.
Las líneas marcadas con \texttt{(*)} son las que lo distinguen del modelo
original, en el que \texttt{s} y \texttt{makespan} tienen dominio
\texttt{0..horizon}, con \texttt{horizon = sum(durations)}:

\noindent\begin{minipage}{\linewidth}\singlespacing\small
\begin{verbatim}
include "globals.mzn";
int: n_jobs;  int: n_machines;
array[1..n_jobs, 1..n_machines] of int: machines;
array[1..n_jobs, 1..n_machines] of int: durations;
int: lb;  int: ub;                                              % (*)
array[1..n_jobs, 1..n_machines] of var 0..ub: s;                % (*)
var lb..ub: makespan;                                           % (*)

constraint forall(j in 1..n_jobs, o in 1..n_machines-1) (
    s[j, o] + durations[j, o] <= s[j, o+1]);
constraint forall(m in 0..n_machines-1) (
    let { array[1..n_jobs] of var int: ms =
            [s[j, o] | j in 1..n_jobs, o in 1..n_machines
             where machines[j, o] == m];
          array[1..n_jobs] of int: md =
            [durations[j, o] | j in 1..n_jobs, o in 1..n_machines
             where machines[j, o] == m];
    } in disjunctive(ms, md));
constraint forall(j in 1..n_jobs) (
    s[j, n_machines] + durations[j, n_machines] <= makespan);
constraint forall(j in 1..n_jobs, o in 1..n_machines) (        % (*)
    s[j, o] >= sum(k in 1..o-1)(durations[j, k]) /\
    s[j, o] + sum(k in o..n_machines)(durations[j, k]) <= ub);
solve minimize makespan;
\end{verbatim}
\end{minipage}
\vspace{\baselineskip}

Como $lb \le C^{*}_{\max} \le ub$, estas restricciones conservan todas las
soluciones óptimas y descartan únicamente los programas peores que el
heurístico. Con ellas los dominios se reducen de forma pronunciada y la
propagación es más efectiva desde el inicio de la búsqueda. A cambio, la
primera solución que encuentre un \textit{solver} ya debe ser al menos tan
buena como la heurística, lo que en instancias grandes puede ser difícil para
un \textit{solver} sin aprendizaje. Por eso el portafolio conserva siempre el
programa heurístico, de manera que si un \textit{solver} no encuentra
solución, el resultado es la $ub$.

\subsection{Modelo con estrategia de búsqueda}
\label{sec:modelo_busqueda}

Otro aporte de este trabajo es obtener del mismo \textit{solver} una
configuración adicional del portafolio (Chuffed V2) mediante una anotación de
búsqueda estándar de MiniZinc, la
cual fija en cada paso la operación con menor inicio posible en su valor
mínimo. Para ello reemplaza la última línea por

{\singlespacing\small
\begin{verbatim}
solve :: int_search([s[j, o] | j in 1..n_jobs, o in 1..n_machines]
                    ++ [makespan], smallest, indomain_min) minimize makespan;
\end{verbatim}
}

\noindent y se ejecuta sin la opción \texttt{-f}, con lo cual el
\textit{solver} aplica esta heurística de ramificación. La búsqueda libre de
Chuffed es muy eficaz para probar la optimalidad en instancias pequeñas y
medianas, pero en las grandes puede tardar en encontrar una primera solución
mejor que la heurística. La estrategia explícita favorece encontrar buenas
soluciones pronto en instancias grandes, a cambio de probar menos óptimos, por
lo que ambas variantes son candidatas naturales a complementarse dentro del
portafolio (Sección~\ref{sec:composicion}).

%------------------------------------------------------------------------------
\section{Formulación del problema de selección y del \textit{schedule}}
\label{sec:formulacion_schedule}
%------------------------------------------------------------------------------

\subsection{Portafolio, instancias y presupuesto de tiempo}
Sea $\Pi = \{a_1, \dots, a_q\}$ el portafolio, $\mathcal{I}$ el conjunto de
instancias y $T$ el presupuesto de tiempo por instancia (300 segundos). Cada
elemento del portafolio es una \textit{configuración}, entendida como un
\textit{solver} junto con el modelo que resuelve. En este trabajo, $\Pi$ tiene
tres configuraciones conformadas por Chuffed y Gecode con el modelo con cotas,
y por Chuffed con el modelo con estrategia de búsqueda (en adelante, Chuffed
V2). Para cada configuración $a \in \Pi$ e instancia $i \in \mathcal{I}$, la
base de conocimiento registra el tiempo $t(a, i)$ que tarda $a$ en probar el
óptimo de $i$ (con $t(a,i) = \infty$ si no lo prueba dentro de $T$) y el mejor
\textit{makespan} $C(a, i, \tau)$ encontrado por $a$ en $i$ tras $\tau$
segundos de ejecución. La función $C(a, i, \tau)$ se obtiene de la
\textit{trayectoria} de la ejecución, es decir, la lista de soluciones
mejoradas que reporta el \textit{solver}, cada una con el instante en que se
encontró. Si $a$ todavía no ha encontrado ninguna solución en el instante
$\tau$, se toma $C(a, i, \tau) = ub(i)$, porque el portafolio dispone siempre
del programa heurístico. Se dice que $a$ resuelve $i$ en un tiempo $\tau$ si
$t(a, i) \le \tau$. Cada instancia se representa además por su vector de
características $f(i) \in \mathbb{R}^{d}$ (Sección~\ref{sec:features}).

Las instancias con $lb(i) = ub(i)$ quedan resueltas por las cotas antes de
ejecutar cualquier configuración. Las comparaciones entre métodos se hacen
sobre las instancias restantes.

\subsection{Selección de un único algoritmo}
En la estrategia \textit{Winner Takes All}, un selector
$S: \mathbb{R}^{d} \rightarrow \Pi$ elige una configuración a partir de las
características de la instancia, y esa configuración se ejecuta durante todo
el presupuesto. El resultado sobre $i$ es el de $S(f(i))$ con tiempo $T$, por
lo cual la instancia se resuelve si $t(S(f(i)), i) \le T$, y la calidad
obtenida es $C(S(f(i)), i, T)$.

Con esta notación se definen las dos referencias de la
Sección~\ref{sec:metricas_portafolios}. El \textbf{SBS} es el selector
constante que elige siempre la configuración con mejor desempeño medio en
$\mathcal{I}$. El \textbf{VBS} es el selector ideal que, para cada instancia,
elige la configuración con mejor desempeño en ella, definida bajo el criterio
de calidad de este trabajo como $\arg\min_{a \in \Pi} C(a, i, T)$.

\subsection{\textit{Schedule} de ejecución}
Un \textit{schedule} de ejecución para la instancia $i$ es una secuencia
\begin{equation}
  \sigma(i) = \bigl[ (a_{(1)}, \tau_1), (a_{(2)}, \tau_2), \dots,
  (a_{(r)}, \tau_r) \bigr],
  \qquad \sum_{l=1}^{r} \tau_l \le T,
  \label{eq:schedule}
\end{equation}
donde $a_{(l)} \in \Pi$ son configuraciones, que pueden repetirse en
ejecuciones independientes, y $\tau_l > 0$ el
tiempo asignado a cada una. Las configuraciones se ejecutan en el orden
indicado, cada una durante su tiempo y de forma independiente, sin transmitir
la mejor solución de una a la siguiente. El resultado del \textit{schedule} es
\begin{equation}
  C(\sigma, i) = \min_{l = 1, \dots, r} C\bigl(a_{(l)}, i, \tau_l\bigr).
  \label{eq:resultado_schedule}
\end{equation}
La instancia se resuelve si alguna configuración la resuelve en su
intervalo, es decir, si existe $l$ tal que $t(a_{(l)}, i) \le \tau_l$, y el
tiempo empleado es la suma de los intervalos anteriores más
$t(a_{(l)}, i)$. Gracias a las trayectorias, la
ecuación~\eqref{eq:resultado_schedule} permite evaluar cualquier
\textit{schedule} sin volver a ejecutar los \textit{solvers}.

Construir un \textit{schedule} implica decidir qué configuraciones incluir
(el subconjunto $\{a_{(1)}, \allowbreak \dots, \allowbreak a_{(r)}\} \subseteq \Pi$), cuánto tiempo
darle a cada una (los valores $\tau_l$) y en qué orden ejecutarlas, que son
las tres decisiones de la pregunta de investigación. La estrategia
\textit{Winner Takes All} es el caso particular $r = 1$ y $\tau_1 = T$. En
este trabajo se comparan varias formas de tomar estas decisiones, que abarcan
un \textit{schedule} fijo para todas las instancias, el enfoque de SUNNY
\citep{amadini2014sunny}, basado en las instancias más parecidas, y un
\textit{schedule} construido con modelos de aprendizaje automático. Su
implementación se describe en el Capítulo~\ref{cap:portafolio}.

%------------------------------------------------------------------------------
\section{Metodología de trabajo}
\label{sec:metodologia}
%------------------------------------------------------------------------------
% Fuente: Anteproyecto, Cap. 7.1.

El desarrollo del proyecto se estructuró en cinco fases, cuyo desglose de
actividades se recoge en la Tabla~\ref{tab:fases}.

\begin{table}[H]
  \centering
  \small
  \singlespacing
  \caption{Fases y actividades del trabajo.}
  \label{tab:fases}
  \begin{tabularx}{\textwidth}{>{\raggedright\arraybackslash}p{0.33\textwidth}
                               >{\raggedright\arraybackslash}X}
    \toprule
    \textbf{Fase} & \textbf{Actividades} \\
    \midrule
    1. Diseñar e implementar un portafolio de algoritmos de optimización para
    el JSSP
      & A1. Identificación y selección de los algoritmos a incluir en el
          portafolio.\newline
        A2. Diseño de las adaptaciones necesarias de los algoritmos para el
          JSSP.\newline
        A3. Implementación de los algoritmos seleccionados.\newline
        A4. Validación del correcto funcionamiento de los algoritmos
          implementados. \\
    \addlinespace
    2. Desarrollar un sistema de caracterización de instancias del problema
      & A5. Identificación de las características relevantes para describir
          instancias del JSSP.\newline
        A6. Diseño del sistema de extracción automática de
          características.\newline
        A7. Implementación del sistema de extracción automática de
          características. \\
    \addlinespace
    3. Construir el conjunto de datos para el entrenamiento del sistema de
    selección
      & A8. Selección de instancias representativas para conformar el conjunto
          de datos.\newline
        A9. Ejecución de los algoritmos sobre las instancias
          seleccionadas.\newline
        A10. Preprocesamiento y preparación del conjunto de datos final. \\
    \addlinespace
    4. Desarrollar el modelo de selección automática de algoritmos
      & A11. Diseño del \textit{pipeline} de \textit{machine learning} para la
          selección de algoritmos.\newline
        A12. Entrenamiento y optimización de los modelos de
          predicción.\newline
        A13. Selección del modelo final con base en métricas de desempeño. \\
    \addlinespace
    5. Implementar el sistema integrado de portafolios de algoritmos
      & A14. Diseño de la arquitectura del sistema integrado.\newline
        A15. Implementación del sistema completo.\newline
        A16. Evaluación experimental y análisis de resultados. \\
    \bottomrule
  \end{tabularx}
\end{table}

Durante el desarrollo, y con la aprobación de la dirección del trabajo, se
hicieron tres ajustes que conservan intactos los objetivos trazados:

\begin{itemize}
  \item En las actividades A2 y A9, el análisis de la primera campaña mostró
    que el horizonte trivial limitaba a Chuffed en las instancias grandes. Se
    incorporaron las cotas de la Sección~\ref{sec:cotas} y se repitió la
    campaña con el modelo acotado, donde la heurística de Giffler y Thompson
    calcula la cota superior y entrega una solución de respaldo.
  \item En las actividades A1 y A2, el portafolio se compuso de tres
    configuraciones basadas en Chuffed y Gecode con el modelo con cotas y en
    Chuffed con estrategia de búsqueda, siguiendo la idea de portafolios de
    configuraciones como Hydra \citep{xu2010hydra}. CP-SAT se empleó como
    referencia externa, dado que domina a los otros dos en el modelo original y
    acapararía la selección (Sección~\ref{sec:composicion}).
  \item En las actividades A11 a A13, la construcción del \textit{schedule} se
    apoyó en un esquema híbrido compuesto por un clasificador de optimalidad
    probada, un regresor del tiempo condicionado a probarla y un clasificador
    de la mejor configuración (Sección~\ref{sec:modelos_prediccion}).
\end{itemize}
